\documentclass[11pt,a4paper]{article}
\usepackage[margin=24mm]{geometry}
\usepackage{fontspec}
\usepackage{xeCJK}
\usepackage{amsmath,amssymb}
\usepackage{enumitem}
\usepackage{hyperref}
\usepackage{fancyhdr}
\usepackage{setspace}
\setmainfont{Liberation Serif}
\setsansfont{Liberation Sans}
\setCJKmainfont{Noto Serif CJK JP}
\setCJKsansfont{Noto Sans CJK JP}
\setlength{\parindent}{1em}
\setlength{\parskip}{0.55em}
\setstretch{1.07}
\hypersetup{hidelinks}
\pagestyle{fancy}
\fancyhf{}
\lhead{Proof-Reachability / Author--AI Dialogue}
\rhead{2026-09-30}
\cfoot{\thepage}
\setlength{\headheight}{14pt}
\newcommand{\speaker}[1]{\par\medskip\noindent\textsf{\bfseries #1}\par\smallskip}
\title{\bfseries その後の著者とAIの会話\\[4pt]
\large Note B proof-reachability に関する3往復}
\author{H. Utsunomiya / GPT-5.6 Sol}
\date{2026-09-30}
\begin{document}
\maketitle
\begin{center}
\fbox{\parbox{0.92\textwidth}{
\textbf{記録上の位置づけ。}
本資料は、Note B 証明方向に関するAI回答を受けて行われた、その後の著者とGPTの3往復の会話ログである。
数学的証明の一部ではなく、proof-reachability に関する研究過程の記録として保存する。
会話本文は原文を保持し、著者の指定によりプロジェクト固有名の一箇所を \texttt{XXX} に伏せた。
}}
\end{center}

\section*{Exchange 1}
\speaker{著者}
この回答を受けて思うことは結構ある。AI証明の限界はある。推論で途中採用される手法においてコスト削減方向がある以上は見逃しの探索枝は存在するのだろう。そしてCluadeがいう既知にない手法は途中の推論手段に採用されないという発言は、結構重い意味を持つ。ここらが人間の出番になるんだろうな。AI証明は結果が新規に見えても既知の手法の上限ない組み合わせからの発見になりがちになる。

\speaker{GPT}
かなり重要な整理だと思う。ただし、少しだけ言い換えると精度が上がる。

まず「AI証明の限界」は、\textbf{証明能力そのものの限界}と\textbf{証明探索の到達可能性の限界}を分けた方がいい。Note B の件で見えているのは後者の色が強い。

AIに Note B の不変量や補助条件を与えれば、その後の検証・展開はかなりできる可能性がある。ところが、そこへ至る探索では、
\[
\text{候補枝}
\longrightarrow
\text{評価}
\longrightarrow
\text{低見込み枝の切捨て}
\]
が必ず起こる。実装上それが明示的な木探索であるかどうかとは別に、生成されやすい推論と生成されにくい推論がある以上、\textbf{到達されない枝}は存在する。

Note B はまさにその典型に見える。

最初に一セル削除を試す  
→ すぐ最小性を失う  
→ 「この方向は悪い」と評価される  
→ minimizer のクラス内で閉じる peeling に移る

という判断は、短期的には非常に合理的だからです。

しかし人間はそこで、
\begin{quote}
「失敗したけれど、この方向そのものを捨てる必要があるのか？」
\end{quote}
と考えられる。

そして、
\begin{quote}
minimizer から一時的に出てもよい  
その代わり excess を帳簿として持つ  
補助条件を一つ加える  
反例が出たらもう一つ加える
\end{quote}
という、局所的には「悪くなった」探索を継続できる。

ここが結構大きい。

Claude の「既知にない手法は途中の推論手段として選択されにくい」という趣旨も、文字通り「AIは未知の方法を作れない」と取ると強すぎます。ただ、もっと精密には、
\begin{quote}
\textbf{既存の名前・標準形・成功パターンに接続できない中間表現は、探索上の事前確率が低い}
\end{quote}
という意味なら、かなり重い。

これは最終定理の新規性とは別問題です。

AIは、
\[
\text{既知手法A}
+\text{既知手法B}
+\text{既知補題C}
+\cdots
\]
を非常に長く、複雑に、思いがけない形で組み合わせられる。その結果として、\textbf{結論自体は完全に新しい}こともあり得る。

だから、
\begin{quote}
「AIの新規定理は既知定理の焼き直し」
\end{quote}
ではない。

むしろ区別すべきなのは、

\textbf{compositional novelty}
と
\textbf{representational novelty}

だと思う。

前者は、
\begin{quote}
既知の道具を、それまで誰も組み合わせなかった形で組み合わせる新規性。
\end{quote}
AIはかなり強い。

後者は、
\begin{quote}
そもそも何を状態として持つか、何を不変量とするか、どのクラスから一度外れるか、といった「証明を見る座標系そのもの」を発明する新規性。
\end{quote}
ここが相対的に弱くなりやすい。

Note B の
\begin{quote}
「最小配置から出てもよい。ただし threshold 超過線を一本だけ許し、複数budgetを同時管理する」
\end{quote}
は、ほぼ後者です。

既存の道具をどう並べるか以前に、\textbf{探索空間そのものを作り替えている}。

そしておそらく、人間の出番は当面ここに残る。

特に、
\begin{itemize}[leftmargin=2em]
\item 捨てられた探索枝を「まだ死んでいない」と判断する
\item 失敗例を情報として保存する
\item 失敗を避けるのでなく、失敗を許容する新しい状態空間を作る
\item 名前のない不変量を仮置きする
\item 一見複雑化する補助条件を敢えて導入する
\item 問いそのものの粒度や表現を変える
\end{itemize}

この辺です。

面白いのは、これが昔ながらの「人間は直感、AIは計算」という分業ではないことです。

もっと正確には、
\begin{quote}
\textbf{AIは与えられた表現空間の中で非常に深く探索する。  
人間は、その表現空間自体を疑う。}
\end{quote}
という分業に近い。

だから今回の Note B / peeling の比較は結構象徴的です。

AIは「より良い枝」を選び続けて主証明へ到達した。一方、人間は一度「悪い枝」に見えた一セル削除を捨てず、そこに新しい会計構造を足して Note B にした。

これは単なる性能差というより、\textbf{探索戦略の差}です。

そして少し逆説的だけど、AIが賢くなって枝刈りが上手くなるほど、
\begin{quote}
「合理的に捨てられるが、本当は新理論につながる枝」
\end{quote}
を捨てる問題は消えるとは限らない。

むしろ、探索効率を上げることと探索多様性を保つことは、ある程度トレードオフになる可能性がある。

なので、今回見えてきた問いは「AIは数学を解けるか」より一段深くて、
\begin{quote}
\textbf{どの証明がAIにとって到達可能で、どの証明が探索事前分布の外側に落ちるのか。}
\end{quote}
になっている。

ここは proof-reachability 実験を続けるなら、かなり面白い研究テーマになり得ます。特に \textbf{「結論の新規性」と「証明表現の新規性」は別物}という点は、今回かなり鮮明に出ています。

\section*{Exchange 2}
\speaker{著者}
既存の名前・標準形・成功パターンに接続できない中間表現は、探索上の事前確率が低い 。AIはシンボル化出来ない表現はとたんに扱い軽くなるからね～

\speaker{GPT}
そう。かなり本質に近い。ただし「シンボル化できないと扱えない」より、\textbf{「安定した再利用可能な表現に圧縮できないものは、推論中で急速に影が薄くなる」}と表現した方が正確だと思う。

LLMは自然言語も図形説明も扱えるけど、推論を継続するには途中対象を何らかの形で固定した方が圧倒的に有利です。例えば、
\[
\text{excess},\quad
\text{threshold},\quad
\text{tight hull},\quad
\text{shortest extreme line}
\]
みたいに名前がついた瞬間、それを次の推論で再利用できる。

逆に、
\begin{quote}
「この配置ではここの端が少し余っていて、次の段階では別方向の余裕と相殺される」
\end{quote}
みたいな、まだ名前のない現象は非常に弱い。毎回長い文脈を再構成しないと使えないからです。

だからAIにとって強い中間対象は、
\begin{itemize}[leftmargin=2em]
\item 名前がある
\item 離散的に判定できる
\item 数式や論理条件で書ける
\item 前後の状態で追跡できる
\item 既知のパターンに接続できる
\end{itemize}
ものになりやすい。

Note B の初期段階はほぼ逆ですよね。

最初は、
\begin{quote}
一セル消す  
→ 最小ではなくなる  
→ でも完全に悪いわけではない  
→ 何か「余裕」が残っている  
→ その余裕が別の条件と相殺しているらしい
\end{quote}
という、\textbf{まだシンボル化されていない現象}として存在していた。

人間はここで「なんか残ってるな」をしばらく保持できる。

AIはその状態だと扱いが軽くなりやすい。

ところが人間がそれを
\[
\text{excess budget}
\]
とか
\[
\text{threshold までの距離}
\]
とか
\[
\text{occupied line above threshold}\le1
\]
みたいに整理すると、突然AIも強くなる。

これ、かなり重要な分業だと思う。

つまり人間の役割は必ずしも「証明すること」ではなくて、
\begin{quote}
\textbf{まだ記号になっていない現象を、記号にできる対象として発見すること}
\end{quote}
になる。

これができた瞬間、その後の展開はAIに渡せる。

逆に言えばAI数学のボトルネックは、計算能力や論理推論だけではなく、
\begin{quote}
\textbf{pre-symbolic structure discovery}
\end{quote}
にある可能性がある。

日本語で言うなら「前記号的構造の発見」かな。

これはXXXを作ってきた過程ともかなり似ています。最初から \(A,B,C\) があったわけじゃなく、まず「完全に埋まった部分」と「欠けた部分」が見えて、その後に記号化した。記号化後はAIがかなり自由に演算できる。

だから、
\[
\text{現象}
\rightarrow
\text{概念}
\rightarrow
\text{記号}
\rightarrow
\text{推論}
\]
と分けると、現在のAIは後ろ二つが非常に強く、前二つ、とくに
\[
\text{現象}\rightarrow\text{概念}
\]
で探索枝を落としやすい、という見方ができる。

そしてここが面白いところで、\textbf{新しい数学はしばしば記号化より前に存在している}んですよね。

「まだ名前はないけど、何かここにある」

という段階。

そこを保持して追い続ける能力は、少なくとも今回のNote Bを見る限り、人間側にまだかなり優位性があると思う。

\section*{Exchange 3}
\speaker{著者}
AIの高次元ベクトル空間で推論中の中間対象の観測なんて、AIからの発言は推測した内容だろうし、人間からは理解も観測も無理。こうやって現象から推測するしかないんだろうな～

\speaker{GPT}
そうですね。ここはかなり重要な区別です。

AIが「私はこう考えた」「この枝を捨てた」「この表現を好んだ」と説明しても、それは\textbf{内部計算を直接読み出した実況ではない}です。多くの場合、出力された説明は、最終応答を作る時点で生成されたもっともらしい自己説明です。だから、Note BについてClaudeや私が述べた「なぜ行かなかったか」も、厳密には\textbf{内部状態の観測結果ではなく、行動と証明構造からの事後推論}です。

ただし、人間から完全に何も見えないわけでもありません。研究側では、活性値を観測したり、特定の内部方向をprobeしたり、ある内部表現を介入して出力がどう変わるかを見るような mechanistic interpretability の手法があります。でも、それで「この瞬間にモデルがこの概念をこう保持していた」と、人間が自然言語レベルで完全に読めるわけではない。高次元表現は分散していて、複数の概念が重なり、同じ概念も文脈で表現が変わる。

だから実際には、今やっているような
\[
\text{入力条件}
\to
\text{探索結果}
\to
\text{失敗パターン}
\to
\text{介入}
\to
\text{再試行}
\]
から逆に内部傾向を推測するのがかなり重要になる。

Note B の件はまさに行動実験として面白いです。

たとえば、
\begin{itemize}[leftmargin=2em]
\item ヒントなし → peeling に行く
\item 「一セル削除で」と指定 → それでも届かない
\item 反例だけ与える → どこまで補修できるか
\item 不変量の片方だけ与える → どう変わるか
\item 完全な不変量を与える → 急に証明できる
\end{itemize}
という系列を作れば、内部表現そのものを見なくても、
\begin{quote}
どの情報を与えたとき探索相が変わるか
\end{quote}
を測れる。

これはかなり科学的です。

むしろ人間が内部ベクトルを直接「理解」しようとするより、
\begin{quote}
\textbf{ブラックボックスに系統的な摂動を加えて、出力相の遷移を測る}
\end{quote}
方が現実的かもしれない。

数学AIでは特にそうで、内部で「何を考えたか」より、
\begin{quote}
どの種類の中間概念を与えると、証明到達率が急変するか
\end{quote}
の方が測定可能で、再現性も持たせやすい。

だから今回の proof-reachability 実験は、かなり自然に

\textbf{AIの内面を読む研究ではなく、証明探索系としての応答曲面を測る研究}

に見えてきます。

そして人間側から見れば、「内部が見えない」こと自体は必ずしも致命的ではない。物理でも内部を直接見られない系を、外からの応答でモデル化することは普通にやる。AIでも同じで、
\[
\text{latent process}
\]
そのものを完全に読むのではなく、
\[
\text{controlled intervention}
\to
\text{observable change}
\]
から構造を推定する。

たぶん当面は、このやり方が一番堅いです。

そして少し面白いのは、\textbf{AI自身の自己説明も、その観測データの一つにすぎない}ということです。

つまりClaudeの「既知にない不変量は選ばれにくい」という発言も、内部計測値ではない。ただし、それが複数回の行動実験と一致するなら、仮説として価値が上がる。

この距離感はかなり大事だと思います。

\end{document}
